% ========================================
% 智慧水利平台架构与开发教材LaTeX配置
% 基于通用教材模板，针对本教材优化
% 注意：本文件不包含documentclass，仅包含包引入和命令定义
% ========================================

% ========================================
% 基础包引入
% ========================================

% 数学支持
\usepackage{amsmath,amssymb,amsthm}

% 图形和表格
\usepackage{graphicx}
\usepackage{booktabs}
\usepackage{array}
\usepackage{longtable}

% 色彩支持
\usepackage{xcolor}

% 页面设置
\usepackage{geometry}
\geometry{
    top=2.5cm,
    bottom=2.5cm,
    left=2.5cm,
    right=2cm,
    bindingoffset=0.5cm
}

% 中文字体设置
\usepackage{fontspec}
\setmainfont{Times New Roman}
\setsansfont{Arial}
\setmonofont{Courier New}

% 设置中文字体
\setCJKmainfont{SimSun}
\setCJKsansfont{SimHei} 
\setCJKmonofont{FangSong}

% 行距设置
\usepackage{setspace}
\onehalfspacing

% 目录和引用
\usepackage{hyperref}
\hypersetup{
    colorlinks=true,
    linkcolor=blue,
    filecolor=magenta,
    urlcolor=cyan,
    citecolor=red,
    bookmarks=true,
    bookmarksopen=true,
    bookmarksnumbered=true,
    pdfstartview=FitH
}

% 列表环境
\usepackage{enumitem}

% 彩色框环境
\usepackage{tcolorbox}
\tcbuselibrary{breakable,skins}

% pifont包用于特殊符号
\usepackage{pifont}

% 代码高亮
\usepackage{listings}
\lstset{
    basicstyle=\ttfamily\small,
    keywordstyle=\color{blue},
    commentstyle=\color{green!60!black},
    stringstyle=\color{red},
    backgroundcolor=\color{gray!10},
    frame=single,
    breaklines=true,
    showstringspaces=false,
    numbers=left,
    numberstyle=\tiny\color{gray}
}

% 定理环境
\newtheorem{theorem}{定理}[chapter]
\newtheorem{definition}[theorem]{定义}
\newtheorem{example}[theorem]{例}

% ========================================
% 颜色定义
% ========================================
\definecolor{primaryblue}{RGB}{41,128,185}
\definecolor{secondarygreen}{RGB}{39,174,96}
\definecolor{warningorange}{RGB}{230,126,34}
\definecolor{dangerred}{RGB}{231,76,60}
\definecolor{lightgray}{RGB}{236,240,241}

% ========================================
% 教育框架命令定义
% ========================================

% 学习目标环境 - 修复版本
\newenvironment{learningobjectives}{%
    \begin{tcolorbox}[
        colback=primaryblue!10,
        colframe=primaryblue,
        title=学习目标,
        fonttitle=\bfseries,
        left=0.5cm,
        right=0.5cm,
        top=0.3cm,
        bottom=0.3cm,
        breakable
    ]%
}{%
    \end{tcolorbox}%
}

% 重点知识 - 修复版本
\newcommand{\keypoints}[1]{%
    \begin{tcolorbox}[
        colback=secondarygreen!10,
        colframe=secondarygreen,
        title=重点知识,
        fonttitle=\bfseries,
        left=0.5cm,
        right=0.5cm,
        top=0.3cm,
        bottom=0.3cm,
        breakable
    ]%
    #1%
    \end{tcolorbox}%
}

% 章节小结 - 修复版本
\newcommand{\chaptersummary}[1]{%
    \begin{tcolorbox}[
        colback=lightgray,
        colframe=gray,
        title=本章小结,
        fonttitle=\bfseries,
        left=0.5cm,
        right=0.5cm,
        top=0.3cm,
        bottom=0.3cm,
        breakable
    ]%
    #1%
    \end{tcolorbox}%
}

% 练习题环境
\newenvironment{exercises}[1][思考题与练习]{%
    \section*{#1}%
    \addcontentsline{toc}{section}{#1}%
}{%
}

% 其他常用命令
\newcommand{\important}[1]{\textbf{\color{dangerred}#1}}
\newcommand{\note}[1]{\textit{\color{primaryblue}#1}}
\newcommand{\warning}[1]{\textbf{\color{warningorange}#1}}

% 页眉页脚设置
\usepackage{fancyhdr}
\pagestyle{fancy}
\fancyhf{}
\fancyhead[LE,RO]{\thepage}
\fancyhead[LO]{\rightmark}
\fancyhead[RE]{\leftmark}
\renewcommand{\headrulewidth}{0.4pt}

% 章节标题格式
\ctexset{
    part/format = \centering\Huge\bfseries,
    chapter/format = \centering\Large\bfseries,
    section/format = \large\bfseries,
    subsection/format = \normalsize\bfseries,
    subsubsection/format = \normalsize
}

% ========================================
% 特殊符号定义（安全定义）
% ========================================
% 使用providecommand避免与已有命令冲突
\providecommand{\crossmark}{\ding{55}}

% 其他有用的命令
\newcommand{\code}[1]{\texttt{#1}}
\newcommand{\file}[1]{\texttt{#1}}
\newcommand{\cmd}[1]{\texttt{#1}}